Notwithstanding the foregoing, in the event any services provided by \lessor, hereunder, or any utilities are interrupted or terminated as a direct result of (i) \lessor's negligence or willful misconduct and not; (ii) curtailment in services imposed by any governmental authority; (iii) failure of the public utilities to furnish necessary services; or (iv) any event or circumstance constituting a force majeure, as indicated in Section~\ref{sec:force-majeure}, (such interruption or termination of services, a ``\dtdef{serviceInterruption}'') and if as a result of such \serviceInterruption, the \premises is untenantable and such \serviceInterruption continues for a period of three (3) consecutive business days after \lessee delivered written notice to \lessor of such \serviceInterruption, then \monthlyRent shall suspend from and including the third (3rd) business day after \lessor's receipt, if such written notice from \lessee, and until such service or utilities are restored to the degree necessary to permit \lessee to reasonably operate its business in the \premises. Furthermore, in the event any such \serviceInterruption continues for a continuous period of more than thirty (30) days and causes the \premises (or a majority of the usable square footage thereof) to be untenantable, then \lessee may, by written notice to \lessor at any time, prior to the resumption of service or utilities to level which allows \lessee to reasonably conduct its normal business operation in the \premises in substantially the manner in which it was used prior to such \serviceInterruption, terminate this \lease. Upon the giving of such notice, this \lease shall terminate on the date set forth in such notice, which date shall be not less than thirty (30) days after the date of such notice.
